\section{Вычислительный юмор как исследовательская проблема}

Юмор давно рассматривается как сложный случай для обработки естественного языка, поскольку успешные шутки зависят не только от грамматической правильности. Они также зависят от несоответствия ожиданиям, общего знания, тайминга и ожиданий аудитории. Ранние обзоры уже описывали юмор как одну из наиболее сложных целей для искусственного интеллекта, отмечая, что реализованные системы в основном ограничивались узкими механизмами каламбуров, а более широкая обработка юмора потребовала бы значительного знания о мире и способности к рассуждению \citep{ritchie2001current}. Два более поздних обзора приходят к сходному выводу с разных сторон. \citet{amin2020survey} показывают, что генерация вычислительного юмора исторически обменивала широту на контролируемость: шаблонные и rule-based системы могли хорошо работать при наличии явного юмористического механизма, но только в сильно ограниченных доменах. Более недавно \citet{Loakman2025} и \citet{lemmens2026computational} утверждают, что эпоха больших языковых моделей расширила покрытие, но не решила центральную проблему, поскольку беглый текст не эквивалентен юмору, который читатели действительно находят смешным.

Дополнительная трудность связана с теоретическим плюрализмом. Теории юмора предлагают несколько повторяющихся классов объяснений---прежде всего подходы, основанные на incongruity, но также традиции superiority и relief,---и при этом ни одна рамка не исчерпывает феномен полностью. Это важно вычислительно, потому что разные типы шуток опираются на разные механизмы. Например, \citet{kao2016computational} формализуют incongruity в каламбурах как неоднозначность плюс отличительность и показывают, что эти переменные предсказывают человеческие оценки смешного. Их работа важна не только тем, что находит применение классической теории юмора, но и тем, что показывает пределы узкой формализации: модель, объясняющая омофонические каламбуры, не начинает автоматически объяснять нарративные шутки, topical jokes или социально ситуированный юмор. Современная theory-driven работа по обнаружению юмора исходит из той же предпосылки: единой всеобъемлющей теории юмора недостаточно, а вычислительные модели обычно усваивают только фрагменты юмора \citep{demarez2024thinc}. Для данной работы следствие является методологическим: юмор следует рассматривать как творческую задачу со множеством потенциальных механизмов, а не как один скалярный атрибут коммуникации.

\section{Эволюция подходов к генерации юмора}

До появления больших языковых моделей системы генерации юмора обычно были узкими по замыслу. Как резюмирует \citet{amin2020survey}, ранние работы часто опирались на вручную написанные шаблоны, лексические ресурсы, фонетические преобразования или жестко ограниченные процедуры генерации. У этих систем было важное преимущество: они кодировали по крайней мере один явный юмористический механизм. Их недостатком была низкая адаптивность. Они генерировали юмор только при тщательно выбранных лингвистических условиях и поэтому плохо масштабировались за пределы каламбуров, устойчивых форм шуток или ограниченных доменов.

Большие языковые модели переворачивают этот trade-off. Они могут генерировать открытый юмористический текст по гораздо более широкому диапазону prompts, но беглость не гарантирует устойчивую юмористическую компетентность. Это видно в недавних крупномасштабных оценках. \citet{Zhang2024HumorAI} вводят мультимодальный benchmark для подписей к карикатурам, построенный более чем на 250 миллионах человеческих оценок для более чем 2.2 миллиона captions, и показывают, что современные методы fine-tuning остаются ограниченными на творческих задачах. Даже сильные системы, такие как GPT-4 и Claude, остаются ниже лучших человеческих участников. Исследования понимания юмора указывают в том же направлении. \citet{xu2024good} показывают, что большие языковые модели демонстрируют то, что они называют lazy pun generation: вместо слияния двух значений в один компактный каламбур модель часто отдельно называет оба связанных слова или значения, поверхностно удовлетворяя задаче. Этот результат методологически важен, поскольку показывает, что простые lexical-success метрики могут переоценивать способность к каламбурам. Оценка должна также отличать структурное качество и оригинальность от механически собранной или скопированной игры слов. \citet{cocchieri2025you} строят устойчивый к contamination benchmark юмора и показывают, что большинство языковых моделей все еще уступают людям в понимании, разрешении и генерации каламбуров. Выходя за пределы каламбуров, \citet{loakman2025comparing} показывают, что современные модели, включая reasoning models, не объясняют надежно все типы шуток, что подчеркивает, насколько большая часть литературы по-прежнему сфокусирована на относительно простых формах юмора.

В совокупности эти исследования показывают, что современные модели лучше воспроизводят поверхностную форму знакомых шуток, чем обобщают юмористические механизмы. Это различие важно для данной работы. Система может хорошо генерировать текст, похожий на шутку, но не быть хорошей в юморе в более сильном смысле, необходимом для открытой генерации.

\section{Проблемы оценки юмора}

Оценка является центральной трудностью вычислительного юмора. Юмор субъективен, но проблема не только в субъективности в разговорном смысле. То, что считается хорошей шуткой, зависит от аудитории, контекста, подачи и оцениваемого измерения. Даже статьи о датасетах подчеркивают, что юмор меняется по теме, времени и получателю, из-за чего трудно выделить одну целевую переменную \citep{weller2020rjokes}. Поэтому современные работы по юмору все чаще избегают трактовки funniness как скалярной метрики.

Наиболее ясный пример дает \citet{Sakabe2025}, оценивающий Oogiri responses по шести измерениям: новизна, ясность, релевантность, intelligence, empathy и общая смешность. Их результаты показывают, что современные языковые модели генерируют ответы между низким и средним уровнем человеческих результатов, но заметно отстают по empathy. Авторы также утверждают, что этот дефицит помогает объяснить, почему model-based evaluation плохо воспроизводит человеческую оценку юмора. Значение этого результата двоякое. Во-первых, он показывает, что качество юмора не сводится к вопросу ``is this funny?'' Во-вторых, он показывает, что человеческие и модельные суждения могут отдавать приоритет разным качествам.

Дополнительная перспектива представлена у \citet{winters2025evaluating}, которые сравнивают человеческих комиков и GPT-4 в живом импровизационном сеттинге, а не в offline text-only benchmark. Их исследование ценно тем, что задает более сильную точку сравнения: и люди, и модель должны быстро отвечать на одни и те же предложения аудитории. Человеческие шутки все еще предпочитаются немного чаще, но разрыв не является подавляющим, а дизайн подчеркивает, насколько сильно подача и ситуативный контекст влияют на оценку юмора. Это полезная коррекция к мышлению только через benchmarks. Текст, хорошо работающий в статической оценке, не обязательно является текстом, который лучше всего работает в реальном комедийном взаимодействии.

Эти результаты также уточняют роль методов LLM-as-a-judge. Структурированные оценочные рамки, такие как у \citet{Hashemi2024}, полезны, потому что заменяют одномерный scoring явными рубриками и калибровкой. Однако эта работа также показывает, что сырые LLM rubric judgments не автоматически хорошо согласуются с человеческими рейтингами. Для генерации юмора это означает, что LLM judges лучше рассматривать как измерительные инструменты с ограниченной областью применимости, а не как прозрачную замену человеческих предпочтений.

\section{Перспектива рассуждения для генерации юмора}

Один из наиболее важных сдвигов в современных исследованиях генерации юмора заключается в переходе от one-shot continuation к явному рассуждению или поэтапному творческому поиску. На уровне prompting \citet{chen2023prompt} моделируют генерацию юмора через пошаговые инструкции комедийного письма для GPT-3. \citet{zhong2024let} идут дальше и утверждают, что стандартный chain-of-thought prompting плохо подходит творческим задачам; они предлагают парадигму Leap-of-Thought для Oogiri-style генерации юмора, где акцент делается на ассоциативных скачках, а не на строго последовательной декомпозиции.

Более явные системы генерации усиливают это направление. \citet{Tikhonov2024} реконструируют написание one-liner jokes как многошаговый процесс и сообщают об улучшениях относительно zero-shot GPT-4 и других baselines в человеческой оценке. \citet{Kim2025} описывают генерацию юмора как комбинацию когнитивных, социальных и творческих навыков. В их исследовании meme-captioning skill-augmented системы существенно превосходят baselines на больших языковых моделях и приближаются к лучшим человеческим captions для целевой аудитории. В мультимодальных условиях \citet{ZhangJ2025} предлагают theory-guided multi-stage reasoning framework, направленный на улучшение и интерпретируемости, и качества юмора.

Эти статьи важны, потому что коллективно отвергают предположение, что юмор лучше всего порождается мгновенным завершением. Вместо этого они моделируют генерацию юмора как процесс поиска по premises, associations, reframings и audience-sensitive choices, аналогично тому, как человеческие эксперты подходят к комедийному письму или другим творческим задачам. Между premise и шуткой может существовать информационный разрыв, который заполняется только во время inference и может быть латентным. В то же время эти подходы разделяют ограничение, напрямую релевантное данной работе: структура рассуждения задается жестко. Число этапов, тип промежуточных представлений и порядок операций фиксируются исследователем. Поэтому такие системы являются сильной демонстрацией полезности рассуждения, но более слабым свидетельством того, что модель может сама научиться, когда и как рассуждать для юмора.

\section{Стандартное постобучение плохо подходит юмору}

Самый прямой путь к улучшению генерации юмора --- supervised fine-tuning на больших корпусах шуток. На практике, однако, такая objective плохо соответствует структуре задачи. Датасеты шуток содержат множество повторяющихся setups, знакомых паттернов wordplay и конвенциональных форм punchline. В результате supervised fine-tuning может улучшать стилистическую правдоподобность, не обязательно улучшая творческое обобщение. Крупные ресурсы юмора, такие как rJokes, ценны масштабом и тематическим разнообразием, но они не были построены как multi-reference generation datasets, где много разных outputs одновременно могут быть правильными \citep{weller2020rjokes}.

Современная литература усиливает это опасение с двух сторон. Во-первых, humor-specific работы показывают, что даже очень большие preference datasets не делают creative alignment простой задачей. \citet{Zhang2024HumorAI} прямо подчеркивают ограничения RLHF- и DPO-style fine-tuning на творческом captioning. Во-вторых, работы по post-training вне юмора делают различие memorization--generalization более резким. \citet{Chu2025} показывают, что reinforcement learning обобщает эффективнее, чем supervised fine-tuning, на unseen textual and visual variants, тогда как supervised fine-tuning склонен запоминать training patterns и испытывать трудности out of distribution. Этот контраст особенно важен для генерации юмора, где weak prompt может иметь много хороших продолжений, а запоминание частых оболочек шуток не равно созданию уместного и неожиданного punchline.

Preference optimization имеет собственные трудности. RLHF предполагает, что reward model может выучить устойчивый proxy для целевого поведения, но суждения о юморе шумны и неоднородны. \citet{Gao2023} в более общем виде показывают, что агрессивная оптимизация несовершенной reward model может ухудшать true performance. Humor-specific evaluation показывает, что риск усиливается в творческих доменах: \citet{Sakabe2025} показывают систематическое расхождение между human and model humor judgments, а \citet{winters2025evaluating} показывают, что качество юмора зависит от аспектов performance, которые трудно захватить offline textual reward signals. По этим причинам ни supervised fine-tuning, ни стандартная preference optimization не решают центральную проблему чисто.

\section{Постобучение, ориентированное на рассуждение}

Недавний успех reasoning-oriented post-training в основном пришел из доменов с объективной верификацией. \citet{Wei2022} показывают, что промежуточное рассуждение может существенно менять поведение модели на задачах, требующих многошагового inference. \citet{Lightman2023} расширяют эту логику от prompting к supervision, показывая, что process supervision превосходит outcome-only supervision на benchmark MATH. \citet{DeepSeek2025} затем демонстрируют, что reinforcement learning может вызывать long-form reasoning без опоры на написанные людьми chains of thought, причем наиболее ясные выигрыши наблюдаются в верифицируемых доменах, таких как математика, программирование и связанные STEM задачи.

Это создает центральную transfer problem для юмора. Юмор также может выигрывать от расширенного рассуждения, но обычно не имеет детерминированного верификатора, способного объявить шутку правильной. Современные verifier-free подходы для доменов вне математики и программирования пытаются закрыть этот разрыв, заменяя явный verifiable reward внутренними вероятностями, назначенными reference answers. \citet{Zhou2025} предлагают VeriFree и показывают, что при single-reference assumption максимизация условной вероятности reference answer при заданных prompt и reasoning trace эквивалентна максимизации exact-match verifier objective в ожидании. \citet{Yu2025} развивают близкий подход RLPR, где reward формулируется как mean likelihood по отдельным токенам для повышения стабильности дисперсии.

Однако для юмора single-reference assumption не является малой технической деталью; это неверная модель данных. Prompts вроде ``write a joke about a frog'' по природе допускают много правдоподобных ответов, и любой датасет будет наблюдать только небольшое и смещенное их подмножество. Иными словами, юмор одновременно является multi-reference и only partially observed задачей. Это означает, что verifier-free reasoning methods концептуально релевантны генерации юмора, но не применимы напрямую без переформулировки reward вокруг множеств референсов, а не одиночных целей.

\section{Исследовательский пробел}

Рассмотренная литература оставляет точный пробел. Исследования генерации юмора все чаще показывают, что staged reasoning, skill decomposition или theory-guided generation могут улучшать outputs, но такие системы в основном опираются на жесткие inference pipelines \citep{chen2023prompt,zhong2024let,Tikhonov2024,Kim2025,ZhangJ2025}. Работы по reasoning-oriented reinforcement learning показывают, что полезные reasoning behaviors можно вызывать обучением, но сильнейшие существующие формулировки предполагают либо объективную верифицируемость, либо один reference answer \citep{Lightman2023,DeepSeek2025,Zhou2025,Yu2025}. В то же время исследования оценки юмора показывают, что funniness является многомерной, контекстно-зависимой и лишь несовершенно захватывается model judges \citep{Zhang2024HumorAI,Sakabe2025,winters2025evaluating,Hashemi2024}.

Данная работа находится на пересечении этих результатов. Она утверждает, что юмор выигрывает от рассуждения, но отвергает необходимость задавать фиксированный творческий сценарий. Она также развивает verifier-free reinforcement learning, но утверждает, что творческие задачи требуют multi-reference formulation, поскольку несколько outputs могут быть одновременно хорошими, а наблюдаемый reference set неполон. Поэтому исследовательский вопрос состоит в том, можно ли адаптировать verifier-free reasoning objectives к partially observed multi-reference humor generation так, чтобы поощрять творческий поиск без схлопывания в memorized training samples.
